\let\negmedspace\undefined
\let\negthickspace\undefined
\documentclass[journal]{IEEEtran}
\usepackage[a5paper, margin=10mm, onecolumn]{geometry}
%\usepackage{lmodern} % Ensure lmodern is loaded for pdflatex
\usepackage{tfrupee} % Include tfrupee package

\setlength{\headheight}{1cm} % Set the height of the header box
\setlength{\headsep}{0mm}     % Set the distance between the header box and the top of the text

\usepackage{gvv-book}
\usepackage{gvv}
\usepackage{cite}
\usepackage{amsmath,amssymb,amsfonts,amsthm}
\usepackage{algorithmic}
\usepackage{graphicx}
\usepackage{textcomp}
\usepackage{xcolor}
\usepackage{txfonts}
\usepackage{listings}
\usepackage{enumitem}
\usepackage{mathtools}
\usepackage{gensymb}
\usepackage{comment}
\usepackage[breaklinks=true]{hyperref}
\usepackage{tkz-euclide} 
\usepackage{listings}
% \usepackage{gvv}                                        
\def\inputGnumericTable{}                                 
\usepackage[latin1]{inputenc}                                
\usepackage{color}                                            
\usepackage{array}                                            
\usepackage{longtable}                                       
\usepackage{calc}                                             
\usepackage{multirow}                                         
\usepackage{hhline}                                           
\usepackage{ifthen}                                           
\usepackage{lscape}
\begin{document}

\bibliographystyle{IEEEtran}
\vspace{3cm}

\title{12.463}
\author{EE25BTECH11033 - Kavin}
% \maketitle
% \newpage
% \bigskip
{\let\newpage\relax\maketitle}

\renewcommand{\thefigure}{\theenumi}
\renewcommand{\thetable}{\theenumi}
\setlength{\intextsep}{10pt} % Space between text and floats
\textbf{Question}:\\
Consider the matrix
\begin{align*}
\vec{A} = \myvec{2 & 1 & 1 \\ 2 & 3 & 4 \\ -1 & -1 & -2}
\end{align*}
whose eigenvalues are $1, -1$ and $3$. Then Trace of $\brak{\vec{A}^3 - 3\vec{A}^2}$ is \rule{1cm}{0.01pt}.
\\    
\bigskip

\textbf{Solution}:\\
The eigenvalues of $\vec{A}$ are given as $\lambda_1 = 1$, $\lambda_2 = -1$, and $\lambda_3 = 3$.\\
\bigskip

We need to find the trace of the matrix $\vec{B} = \vec{A}^3 - 3\vec{A}^2$. This is a polynomial of the matrix $\vec{A}$, where the polynomial is $p(x) = x^3 - 3x^2$.\\
\bigskip

If a matrix $\vec{A}$ has eigenvalues $\lambda_1, \lambda_2, \dots, \lambda_n$, then for any polynomial $p(x)$, the matrix $p(\vec{A})$ will have eigenvalues $p(\lambda_1), p(\lambda_2), \dots, p(\lambda_n)$.\\
\bigskip

Using this property, the eigenvalues of $\vec{B}$ (let's call them $\mu_i$) are $p(\lambda_i)$.

\begin{itemize}
    \item For $\lambda_1 = 1$:
    $$ \mu_1 = (1)^3 - 3(1)^2 = 1 - 3 = -2 $$
    
    \item For $\lambda_2 = -1$:
    $$ \mu_2 = (-1)^3 - 3(-1)^2 = -1 - 3(1) = -1 - 3 = -4 $$
    
    \item For $\lambda_3 = 3$:
    $$ \mu_3 = (3)^3 - 3(3)^2 = 27 - 3(9) = 27 - 27 = 0 $$
\end{itemize}
So, the eigenvalues of the matrix $(\vec{A}^3 - 3\vec{A}^2)$ are -2, -4, and 0.\\
\bigskip


The trace of a matrix is the sum of its eigenvalues.
\begin{align*}
\text{Trace}(\vec{A}^3 - 3\vec{A}^2) &= \mu_1 + \mu_2 + \mu_3 \\
&= (-2) + (-4) + 0 \\
&= -6
\end{align*}

The Trace of $(\vec{A}^3 - 3\vec{A}^2)$ is -6.

\end{document}


